\documentclass[11pt,a4paper]{article}
\usepackage[utf8]{inputenc}
\usepackage[spanish]{babel}
\usepackage{amsmath}
\usepackage{amssymb}
\usepackage{graphicx}
\usepackage{listings}
\usepackage{xcolor}
\usepackage{geometry}
\geometry{
    a4paper,
    total={170mm,257mm},
    left=20mm,
    top=20mm,
}

\definecolor{codegreen}{rgb}{0,0.6,0}
\definecolor{codegray}{rgb}{0.5,0.5,0.5}
\definecolor{codepurple}{rgb}{0.58,0,0.82}
\definecolor{backcolour}{rgb}{0.95,0.95,0.92}

\lstdefinestyle{mystyle}{
    backgroundcolor=\color{backcolour},   
    commentstyle=\color{codegreen},
    keywordstyle=\color{magenta},
    numberstyle=\tiny\color{codegray},
    stringstyle=\color{codepurple},
    basicstyle=\ttfamily\footnotesize,
    breakatwhitespace=false,         
    breaklines=true,                 
    captionpos=b,                    
    keepspaces=true,                 
    numbers=left,                    
    numbersep=5pt,                  
    showspaces=false,                
    showstringspaces=false,
    showtabs=false,                  
    tabsize=2
}
\lstset{style=mystyle}

\title{Proyecto Arbitraje Híbrido en CEX:\\Arquitectura y Matemática}
\author{Documento Maestro}
\date{\today}

\begin{document}

\maketitle
\tableofcontents
\newpage

\section{Fundamentos Teóricos y Modelado Matemático (El Laboratorio)}

Aquí introduciremos los tipos de ineficiencias en los exchanges centralizados (CEX), y las limitaciones lógicas y físicas como latencia y comisiones.

\subsection{Teoría de Grafos Aplicada al Arbitraje Triangular}

El arbitraje triangular busca explotar ineficiencias de precio entre tres pares de divisas dentro de un mismo exchange para evitar la latencia de red entre servidores distintos. Un ejemplo clásico es: USDT $\to$ BTC $\to$ ETH $\to$ USDT.

Para que exista una ganancia real, el producto de las tasas de cambio (considerando las comisiones) tiene que ser estrictamente mayor a 1.
Si $R$ es la tasa de cambio y $f$ es la comisión (fee), queremos que:
\begin{equation}
(R_1 \times (1 - f)) \times (R_2 \times (1 - f)) \times (R_3 \times (1 - f)) > 1
\end{equation}

Trabajar con multiplicaciones en algoritmos de búsqueda puede ser ineficiente e inducir problemas de precisión flotante. Por eso aplicamos logaritmo natural a ambos lados. Sabiendo que $\log(A \times B) = \log(A) + \log(B)$, la inecuación se transforma en una suma:
\begin{equation}
\log(R_1(1-f)) + \log(R_2(1-f)) + \log(R_3(1-f)) > 0
\end{equation}

Para aplicar algoritmos estándar de grafos que buscan el camino "más corto" o "mínimo", como el de Bellman-Ford, multiplicamos todo por $-1$, invirtiendo la desigualdad:
\begin{equation}
-\log(R_1(1-f)) - \log(R_2(1-f)) - \log(R_3(1-f)) < 0
\end{equation}

\textbf{Conclusión Teórica:} 
Es posible modelar el mercado como un grafo dirigido $G = (V, E)$, donde los vértices ($V$) son las monedas y las aristas ($E$) representan los mercados disponibles. El peso de cada arista es $w = -\log(R \times (1 - f))$.
Si el algoritmo de \textbf{Bellman-Ford} encuentra un ciclo con peso total negativo en este grafo, podemos afirmar que existe una oportunidad matemática de arbitraje rentable.

\vspace{0.5cm}
Implementación base en Python para evaluar el modelo sobre el grafo:

\begin{lstlisting}[language=Python]
import math

class ArbitrajeTriangular:
    def __init__(self, fee=0.001): 
        # Asumimos 0.1% fee estandar (ej. Bybit/OKX Taker)
        self.fee = fee
        self.grafo = []
        self.monedas = []

    def agregar_mercado(self, base, quote, precio_bid, precio_ask):
        if base not in self.monedas: self.monedas.append(base)
        if quote not in self.monedas: self.monedas.append(quote)

        # Arista: Vender Base para obtener Quote (Peso = -log(bid * (1 - fee)))
        if precio_bid > 0:
            peso_venta = -math.log(precio_bid * (1 - self.fee))
            self.grafo.append((base, quote, peso_venta))

        # Arista: Comprar Base usando Quote (Peso = -log((1 / ask) * (1 - fee)))
        if precio_ask > 0:
            peso_compra = -math.log((1.0 / precio_ask) * (1 - self.fee))
            self.grafo.append((quote, base, peso_compra))

    def bellman_ford(self, nodo_origen):
        distancias = {nodo: float('inf') for nodo in self.monedas}
        predecesores = {nodo: None for nodo in self.monedas}
        distancias[nodo_origen] = 0

        # Relajacion de aristas (V - 1 veces)
        for _ in range(len(self.monedas) - 1):
            for u, v, w in self.grafo:
                if distancias[u] + w < distancias[v]:
                    distancias[v] = distancias[u] + w
                    predecesores[v] = u

        # Deteccion de ciclos negativos
        ciclos = []
        for u, v, w in self.grafo:
            if distancias[u] + w < distancias[v] - 1e-9:
                ciclo = []
                nodo_actual = v
                for _ in range(len(self.monedas)):
                    nodo_actual = predecesores[nodo_actual]
                
                nodo_inicio_ciclo = nodo_actual
                ciclo.append(nodo_inicio_ciclo)
                nodo_actual = predecesores[nodo_inicio_ciclo]
                
                while nodo_actual != nodo_inicio_ciclo:
                    ciclo.append(nodo_actual)
                    nodo_actual = predecesores[nodo_actual]
                ciclo.append(nodo_inicio_ciclo)
                
                ciclo.reverse()
                if ciclo not in ciclos:
                    ciclos.append(ciclo)

        return ciclos
\end{lstlisting}

\section{Arquitectura del Sistema (El Stack H\'ibrido)}

El sistema est\'a dividido en dos grandes subsistemas desacoplados que se comunican v\'ia \textbf{Redis}:

\begin{itemize}
    \item \textbf{El M\'usculo} (\texttt{ingesta\_websockets.py}): Motor de ingest\'{\i}on. Conecta WebSockets de Bybit v\'ia CCXT Pro y vuelca el mejor \textit{bid/ask} de cada par a Redis en tiempo real con latencia sub-milisegundo.
    \item \textbf{El C\'alculo} (\texttt{detector\_arbitraje.py}): Motor de an\'alisis. Lee el grafo desde Redis cada 50ms y ejecuta Bellman-Ford. Si encuentra un ciclo negativo, dispara la alerta de oportunidad.
    \item \textbf{El Cerebro} (\texttt{main.py}): El orquestador. Lanza ambos motores en paralelo usando \texttt{asyncio.gather}.
\end{itemize}

\section{Motor de Ejecuci\'on y Conectividad (El M\'usculo)}

Para encontrar estas oportunidades rentables es imperativo recolectar la microestructura del mercado (Limit Order Book) de forma casi instantánea. 

\subsection{Conexión en Tiempo Real con Exchanges (WebSockets)}

Se diseñó un módulo asíncrono utilizando \textbf{CCXT Pro} (basado en Asyncio), el cual mantiene conexiones WebSocket persistentes contra exchanges de primera línea (ej. Bybit). En lugar de usar la ineficiente arquitectura REST, el WebSocket solo devuelve \textit{deltas} o actualizaciones nuevas del mercado, permitiendo latencias asimétricas mínimas.

Para desacoplar el motor de consumo (la red) del motor de cálculo matemático (Bellman-Ford), se integró una base de datos en memoria ultrarrápida bidireccional: \textbf{Redis}.
A medida que el socket recibe las actualizaciones de precios para los tres vértices que conforman el triángulo ($BTC\to USDT$, $ETH\to BTC$, $ETH\to USDT$), el mejor bid y ask es inyectado inmediatamente mediante un \textit{Hash} `O(1)` en Redis.

\begin{lstlisting}[language=Python]
async def stream_orderbook(exchange, r_client, symbol):
    # CCXT Pro usa WebSocket nativo
    orderbook = await exchange.watch_order_book(symbol)
    best_bid = orderbook['bids'][0][0]
    best_ask = orderbook['asks'][0][0]
    
    # Push ultrarrapido a Redis
    await r_client.hset(
        f"bybit:{symbol}", 
        mapping={'bid': best_bid, 'ask': best_ask, 'time': exchange.milliseconds()}
    )
\end{lstlisting}

\subsection{Motor de Detecci\'on (Bellman-Ford en Tiempo Real)}

El m\'odulo \texttt{detector\_arbitraje.py} corre un ciclo as\'incrono puro que se ejecuta aproximadamente 20 veces por segundo (cada $50$ms). En cada iteraci\'on:

\begin{enumerate}
    \item Lee los precios actuales de los 3 pares desde Redis mediante \texttt{HGETALL}.
    \item Reconstruye el grafo dirigido $G = (V, E)$ con los pesos logar\'itmicos actualizados: $w = -\log(R \times (1 - f))$.
    \item Ejecuta Bellman-Ford desde cada v\'ertice.
    \item Si se detecta un ciclo negativo, registra la oportunidad con la ruta completa y el timestamp.
\end{enumerate}

\section{Orquestaci\'on e Inteligencia Artificial (El Cerebro)}

\subsection{Interfaz de Control (Telegram)}

Se implement\'o un m\'odulo de alertas nativo sobre la API de Telegram para notificar al operador en tiempo real. El m\'odulo \texttt{telegram\_alertas.py} env\'ia mensajes formateados ante dos eventos cr\'iticos:

\begin{itemize}
    \item \textbf{Oportunidad aprobada:} El ciclo negativo superó la validación del Risk Engine y el capital estimado de ganancia se muestra al instante.
    \item \textbf{Circuit Breaker:} Si se acumulan $N$ errores consecutivos en una ventana de tiempo, el sistema se detiene y notifica al operador para intervención manual.
\end{itemize}

\section{Gesti\'on de Riesgo (Risk Engine)}

El m\'odulo \texttt{risk\_engine.py} act\'ua como guardia antes de que cualquier oportunidad detectada por Bellman-Ford sea ejecutada o notificada. Implementa tres capas de protecci\'on:

\begin{enumerate}
    \item \textbf{Filtro de Spread Neto:} Calcula el spread neto descontando los fees de cada eslab\'on del ciclo. Solo pasan oportunidades donde $\text{spread}_{\text{neto}} > \delta_{\min}$, siendo $\delta_{\min}$ el umbral configurable (por defecto $0.05\%$).
    \item \textbf{Filtro de Liquidez:} Compara el capital m\'aximo configurado contra la profundidad disponible en el \textit{order book}. Si la profundidad es menor a un m\'inimo, la oportunidad se descarta para evitar \textit{slippage} excesivo.
    \item \textbf{Circuit Breaker:} Contador de errores con ventana temporal de $T$ segundos. Si se acumulan $N$ errores consecutivos, el sistema activa el circuit breaker y detiene toda operaci\'on hasta reinicio manual por el operador.
\end{enumerate}

\section{Estrategias Algor\'itmicas de Producci\'on}

\begin{itemize}
    \item \textbf{Arbitraje Triangular de Baja Latencia:} Implementado y funcionando (Secciones I y III).
    \item \textbf{Spot-Perp Basis Arbitrage:} En desarrollo -- fase II.
    \item \textbf{Funding Rate Arbitrage:} En desarrollo -- fase II.
\end{itemize}

\section{Pruebas y Backtesting}

Las pruebas de validaci\'on del sistema se ejecutaron sobre datos artificiales con spreads conocidos (ver Secci\'on I.2.). En la siguiente fase se implementar\'a un backtester con datos tick-by-tick descargados directamente de los exchanges.

\end{document}

